\subsection{Основные теоретические подходы: концепция <<гражданской культуры>> Г. Алмонда и С. Вербы}

Вторым важнейшим этапом в методологии исследования политической культуры стало появление эмпирического и сравнительного подходов, наиболее полное воплощение нашедших в классическом труде американских политологов Г. Алмонда и С. Вербы <<Гражданская культура и стабильность демократии>>. В отличие от нормативного подхода, авторы сместили фокус исследовательского внимания на изучение субъективного измерения политики --- на совокупность психологических ориентаций граждан по отношению к политической системе и ее различным элементам. Под политической культурой нации исследователи понимали специфическое распределение образцов политических ориентаций (познавательных, эмоциональных и оценочных) относительно объектов власти\footnote{Алмонд Г., Верба С. Гражданская культура и стабильность демократии // Полис. Политические исследования. --- 1992. --- \textnumero~4. --- С. 122.}.

На основе масштабного сравнительного исследования пяти стран Г. Алмонд и С. Верба разработали идеальную типологию политических культур, выделив три «чистых» (базовых) типа, различающихся степенью вовлеченности граждан в политический процесс:

1. Патриархальная (парохиальная) политическая культура. Характеризуется практически полным отсутствием интереса граждан к центральным институтам государственной власти. У индивидов отсутствуют ожидания от политической системы, а их ориентации направлены исключительно на локальные, племенные или общинные структуры. В данном типе культуры не существует специализированных политических ролей.

2. Подданническая политическая культура. Для нее характерен высокий уровень осведомленности граждан о функционировании исполнительных органов власти и результатах государственной политики (так называемых «выходах» системы). Однако индивиды воспринимают себя исключительно как объекты управления, демонстрируя пассивность, отсутствие инициативы по отношению к «входам» системы (каналам формирования политической воли) и стремление к безусловному подчинению.

3. Активистская (участническая) политическая культура. Отличается максимальной вовлеченностью населения в политическую жизнь. Граждане демонстрируют высокую информированность и развитые навыки участия как на стадии принятия решений элитами, так и на стадии оценки результатов управления. Индивиды осознают себя полноценными субъектами политического процесса, способными оказывать влияние на институты власти\footnote{Алмонд Г., Верба С. Гражданская культура и стабильность демократии // Полис. Политические исследования. --- 1992. --- \textnumero~4. --- С. 124.}.

Тем не менее главным открытием исследователей стало доказательство того, что ни один из «чистых» типов в изолированном виде не способен обеспечить устойчивое функционирование демократического режима. Стабильность демократии гарантируется возникновением особого смешанного типа, названного авторами <<гражданской культурой>> (*Civic Culture*). 

Гражданская культура представляет собой уникальный баланс, в котором активистские устремления большинства населения уравновешиваются и смягчаются элементами подданнической преданности государству и патриархальной привязанности к традиционным ценностям. В этой модели граждане достаточно активны, чтобы их мнение учитывалось элитами, но одновременно с этим достаточно законопослушны, чтобы обеспечивать управляемость государства и легитимность принимаемых решений. Таким образом, гражданская культура разрешает фундаментальное противоречие демократии между необходимостью поддержания государственной власти и требованием гражданского участия\footnote{Алмонд Г., Верба С. Гражданская культура и стабильность демократии // Полис. Политические исследования. --- 1992. --- \textnumero~4. --- С. 131.}.